\section{Introduction}
\subsection{Motivation}
Transformer-based models \cite{vaswani2017attention} have achieved state-of-the-art results across a wide range of machine learning tasks, including natural language processing (NLP), computer vision, and speech processing. The core innovation of the Transformer is the \textbf{self-attention} mechanism, which allows each element in a sequence to directly interact with every other element, capturing long-range dependencies without the sequential bottleneck of recurrent networks. However, this expressiveness comes at a steep computational price: for a sequence of length $N$ and hidden dimension $d$, self-attention requires computing an $N \times N$ attention matrix, leading to $O(N^2 d)$ time complexity and $O(N^2)$ memory complexity. This quadratic scaling becomes prohibitive for long sequences—for example, processing a document of 10,000 tokens would require 100 million attention entries per layer, quickly exceeding the memory of typical GPUs.

Numerous approaches have been proposed to address this limitation, falling into several categories: sparse attention patterns (e.g., Longformer \cite{beltagy2020longformer}, BigBird \cite{zaheer2020bigbird}), kernel-based linear attention (e.g., Performer \cite{choromanski2020rethinking}), and low-rank approximations (e.g., Linformer \cite{wang2020linformer}). Among these, low-rank methods are particularly appealing because they preserve the full attention mechanism but operate on a compressed representation, often yielding a direct trade-off between accuracy and efficiency via the rank parameter.

\subsection{Limitations of Existing Low-Rank Methods}
The most prominent low-rank attention method, Linformer, projects the keys and values using a fixed random matrix $E \in \mathbb{R}^{N \times r}$ (or $F \in \mathbb{R}^{N \times r}$), independent of the input. While this reduces complexity to $O(N r d)$, the fixed projection may not be optimal for all inputs because it does not adapt to the specific structure of each sequence. In information-theoretic terms, a fixed projection can only capture generic correlations and may discard input-specific features that are crucial for the task.

\subsection{Our Contribution: Data-Dependent SVD Projection}
We propose to replace the fixed random projection with a \textbf{data-dependent projection} derived from the input sequence itself using truncated Singular Value Decomposition (SVD). Specifically, for each input hidden state $H \in \mathbb{R}^{N \times d}$, we compute its top-$r$ left singular vectors $U_r \in \mathbb{R}^{N \times r}$, which form an orthonormal basis that best approximates the row space of $H$ in the Frobenius norm (Eckart–Young theorem). We then project the keys $K$ and values $V$ using $U_r^\top$, effectively compressing the sequence dimension from $N$ to $r$. The attention scores are computed as $Q (U_r^\top K)^\top$, which yields an $N \times r$ matrix instead of $N \times N$.

This approach has several advantages:
\begin{enumerate}
    \item \textbf{Optimal low-rank approximation:} For a given rank $r$, $U_r \Sigma_r V_r^\top$ is the best rank-$r$ approximation to $H$ in both Frobenius and spectral norms.
    \item \textbf{Adaptivity:} The projection changes per input, capturing task-specific patterns.
    \item \textbf{Simplicity:} SVD is a well-understood linear algebra operation with efficient randomized algorithms.
    \item \textbf{Theoretical grounding:} The approximation error is bounded by the $(r+1)$-th singular value, which decays rapidly for many real-world datasets.
\end{enumerate}

\subsection{Paper Outline}
The remainder of this paper is organized as follows. Section 2 provides a comprehensive background on self-attention, SVD, and low-rank matrix approximation. Section 3 presents our SVD-Lin method in detail, including mathematical formulation, complexity analysis, and integration into BERT. Section 4 describes the experimental setup, including GLUE tasks, hyperparameters, and hardware (M1 Mac). Section 5 reports results on accuracy, efficiency, and ablation studies. Section 6 discusses the implications, limitations, and potential extensions. Section 7 concludes the paper. All code and data are publicly available.